% T33_Sub_2Node.tex
% Session: Subtraction 2-Node Construction and SuperBEST v3 Update
% Date: 2026-04-20
% Status: T33 (sub = 2 nodes in F16), SuperBEST v3 (19 nodes / 74.0% savings)

\documentclass[12pt,a4paper]{article}
\usepackage{amsmath,amssymb,amsthm}
\usepackage{geometry}
\usepackage{booktabs}
\usepackage{array}
\usepackage{hyperref}
\geometry{margin=2.5cm}

% ── Theorem environments ────────────────────────────────────────────────────
\newtheorem{theorem}{Theorem}[section]
\newtheorem{lemma}[theorem]{Lemma}
\newtheorem{corollary}[theorem]{Corollary}
\newtheorem{proposition}[theorem]{Proposition}
\theoremstyle{definition}
\newtheorem{definition}[theorem]{Definition}
\newtheorem{remark}[theorem]{Remark}

% ── Operator macros ─────────────────────────────────────────────────────────
\newcommand{\EML}{\operatorname{EML}}
\newcommand{\EDL}{\operatorname{EDL}}
\newcommand{\EXL}{\operatorname{EXL}}
\newcommand{\EAL}{\operatorname{EAL}}
\newcommand{\EMN}{\operatorname{EMN}}
\newcommand{\DEML}{\operatorname{DEML}}
\newcommand{\ELAd}{\operatorname{ELAd}}
\newcommand{\ELSb}{\operatorname{ELSb}}
\newcommand{\LEAd}{\operatorname{LEAd}}
\newcommand{\LEdiv}{\operatorname{LEdiv}}
\newcommand{\LEprod}{\operatorname{LEprod}}
\newcommand{\EPL}{\operatorname{EPL}}
\newcommand{\DEAL}{\operatorname{DEAL}}
\newcommand{\DEXL}{\operatorname{DEXL}}
\newcommand{\DEDL}{\operatorname{DEDL}}
\newcommand{\DEPL}{\operatorname{DEPL}}
\newcommand{\DEMN}{\operatorname{DEMN}}
\newcommand{\LEX}{\operatorname{LEX}}
\newcommand{\sub}{\operatorname{sub}}
\newcommand{\mul}{\operatorname{mul}}

\title{Subtraction in Two Nodes:\\
       T33 and the SuperBEST v3 Table Update to 19 Nodes}
\author{Arturo R.\ Almaguer\\
\small monogate research \quad \texttt{almaguer1986@gmail.com}}
\date{April 2026}

% ════════════════════════════════════════════════════════════════════════════
\begin{document}
\maketitle

% ── Abstract ────────────────────────────────────────────────────────────────
\begin{abstract}
We prove that $\sub(x,y) = x - y$ can be computed by a \emph{two-node} tree
over the extended 16-operator family $\mathcal{F}_{16}$, improving the
previous best of three nodes.  The construction is
\[
  \LEdiv\!\bigl(x,\;\EML(y,\,1)\bigr) \;=\; x - y,
\]
where Node~1 computes $\EML(y,1)=e^y$ and Node~2 computes
$\LEdiv(x,e^y)=\ln(e^x/e^y)=x-y$.
The one-node lower bound is certified by exhaustive search over all 256 one-node
candidate trees in $\mathcal{F}_{16}$, finding zero matches.
A complementary partial-derivative argument shows why no single operator
simultaneously achieves $\partial_x = +1$ and $\partial_y = -1$ over all
$(x,y)$.
The exhaustive two-node search further finds five additional constructions,
confirming that the minimum is exactly two.
Together with the previous T10u result (mul also requires 2 nodes), this
reduces the SuperBEST total from 20 to \textbf{19 nodes} and raises the
compression savings to $1 - 19/73 \approx \mathbf{74.0\%}$.
\end{abstract}

\tableofcontents
\bigskip

% ════════════════════════════════════════════════════════════════════════════
\section{Setup and Definitions}
% ════════════════════════════════════════════════════════════════════════════

\begin{definition}[The 16-operator extended family $\mathcal{F}_{16}$]
\label{def:F16}
The \emph{six-operator library} $\mathcal{F}_6$ consists of:
\[
  \mathcal{F}_6 \;=\; \{\EML,\,\EDL,\,\EXL,\,\EAL,\,\EMN,\,\DEML\},
\]
where every operator has the form $e^{\pm a} \star \ln b$ for arithmetic
$\star \in \{-,+,\times,\div\}$.
The \emph{extended 16-operator family} $\mathcal{F}_{16}$ adds ten further
operators, including notably:
\begin{align*}
  \LEdiv(a,b) &= \ln\!\bigl(e^{a}/b\bigr) = a - \ln b, \\
  \LEprod(a,b) &= \ln\!\bigl(e^{a} \cdot b\bigr) = a + \ln b, \\
  \ELAd(a,b)  &= e^{a} \cdot b, \\
  \ELSb(a,b)  &= e^{a} / b.
\end{align*}
\end{definition}

\begin{definition}[$k$-node tree; terminal set]\label{def:tree}
Fix a terminal set $\mathcal{T} = \{x,\,y,\,0,\,1\}$.
A \emph{$k$-node tree} over operator family $\mathcal{F}$ is a complete
binary tree with exactly $k$ internal nodes labeled by operators in
$\mathcal{F}$; leaves are labeled by elements of $\mathcal{T}$.
For a two-node tree, the two canonical shapes are:
\begin{itemize}
  \item \textbf{Shape~A}: $\mathrm{op}_2\!\bigl(\mathrm{op}_1(a,b),\,c\bigr)$
        \quad (left subtree internal),
  \item \textbf{Shape~B}: $\mathrm{op}_2\!\bigl(a,\,\mathrm{op}_1(b,c)\bigr)$
        \quad (right subtree internal).
\end{itemize}
\end{definition}

\begin{definition}[Witness set for subtraction]\label{def:witness}
We use the \emph{8-point witness set}
\[
  \mathcal{W} \;=\;
  \{(5,3),\;(10,7),\;(1,1),\;(-2,3),\;(0.5,0.5),\;(e,1),\;(2,e),\;(3.5,1.5)\}
\]
with tolerance $\varepsilon=10^{-9}$.  A tree \emph{passes the witness test}
if $|T(x,y)-(x-y)|<\varepsilon$ at every point of $\mathcal{W}$.
\end{definition}

% ════════════════════════════════════════════════════════════════════════════
\section{The One-Node Lower Bound}
% ════════════════════════════════════════════════════════════════════════════

\begin{theorem}[T33a: No 1-node tree computes $x-y$ in $\mathcal{F}_{16}$]
\label{thm:lower1}
There is no 1-node tree over $\mathcal{F}_{16}$ and terminal set
$\mathcal{T}=\{x,y,0,1\}$ that computes $\sub(x,y)=x-y$.
\end{theorem}

\begin{proof}
We give both an exhaustive search and a structural argument.

\smallskip
\noindent\textbf{Part~I: Exhaustive search.}

The 1-node search space has $|\mathcal{F}_{16}|\times|\mathcal{T}|^2 = 16\times 16 = 256$ candidate trees.
The script \texttt{python/scripts/door5\_sub\_2node.py} evaluates each
candidate at $\mathcal{W}$ (Definition~\ref{def:witness}) and finds:
\[
  \texttt{lower\_bound\_1node}\;:\;\texttt{matches\_found} = 0.
\]
No 1-node tree passes the witness test, confirming $\sub \notin \mathcal{F}_{16}^{(1)}$.

\smallskip
\noindent\textbf{Part~II: Structural (partial-derivative) argument.}

For a 1-node tree $T = \mathrm{op}(a,b)$ with $a,b\in\mathcal{T}$ to equal $x-y$,
we need the partial derivatives to satisfy:
\[
  \frac{\partial T}{\partial x} = +1 \quad\text{and}\quad
  \frac{\partial T}{\partial y} = -1 \quad \text{simultaneously for all } (x,y).
\]
We categorize all operators in $\mathcal{F}_{16}$ by their structural form:

\begin{enumerate}
  \item \textbf{Form $e^{\pm a}\star\ln b$} (the $\mathcal{F}_6$ subfamily and its
        negated variants): any output involving $e^{\pm x}$ has
        $\partial/\partial x = \pm e^{\pm x}$, which is not identically $\pm 1$.
        Any output involving $\ln y$ has $\partial/\partial y = 1/y$,
        not identically $-1$.  No constant partial derivatives arise.

  \item \textbf{Form $e^{a}\cdot b$ or $e^{a}/b$} (ELAd, ELSb):
        $\partial/\partial a = e^a \cdot b$ or $e^a/b$, transcendental in $a$.
        Cannot equal $\pm 1$ for all $a$.

  \item \textbf{Form $\ln(e^{a}\pm b)$} (LEAd, LEX):
        $\partial/\partial b = \pm 1/(e^{a}\pm b)$, which is not constant.

  \item \textbf{Form $\ln(e^{a}\cdot b)$ or $\ln(e^{a}/b)$} (\textbf{LEprod, LEdiv}):
        $\LEdiv(x,y) = x - \ln y$, so $\partial/\partial x = 1$ but
        $\partial/\partial y = -1/y \neq -1$ in general.
        $\LEprod(x,y) = x + \ln y$, so $\partial/\partial y = 1/y \neq -1$.
\end{enumerate}

In particular, $\LEdiv(x,y) = x - \ln y$ is the ``closest'' operator,
achieving $\partial_x = 1$ but $\partial_y = -1/y \neq -1$.
No operator achieves both constant partial derivatives simultaneously.
This structural argument corroborates the exhaustive search.
\end{proof}

% ════════════════════════════════════════════════════════════════════════════
\section{T33: The Two-Node Construction}
% ════════════════════════════════════════════════════════════════════════════

\subsection{The $\LEdiv$/$\EML$ construction}

\begin{theorem}[T33: $\sub(x,y) = 2$ nodes in $\mathcal{F}_{16}$]\label{thm:T33}
For all $(x,y)\in\mathbb{R}\times\mathbb{R}_{>0}$,
\[
  \boxed{%
    \sub(x,y) \;=\; \LEdiv\!\bigl(x,\;\EML(y,\,1)\bigr) \;=\; x - y.
  }
\]
This two-node tree over $\mathcal{F}_{16}$ is \emph{optimal}: no one-node
tree in $\mathcal{F}_{16}$ computes $x-y$ (Theorem~\ref{thm:lower1}).
\end{theorem}

\begin{proof}
\noindent\textbf{Step 1: Inner node (Node 1).}
\[
  \EML(y,\,1) \;=\; e^{y} - \ln 1 \;=\; e^{y} - 0 \;=\; e^{y}.
\]

\noindent\textbf{Step 2: Outer node (Node 2).}
\[
  \LEdiv\!\bigl(x,\;e^{y}\bigr)
  \;=\; \ln\!\bigl(e^{x} / e^{y}\bigr)
  \;=\; \ln\!\bigl(e^{x-y}\bigr)
  \;=\; x - y.
\]

\noindent\textbf{Step 3: Node count.}
Node~1 is $\EML(y,1)$; Node~2 is $\LEdiv(x,\,\cdot\,)$.
Total: \textbf{2 nodes}.

\noindent\textbf{Step 4: Numerical verification.}
The script \texttt{python/scripts/door5\_sub\_2node.py} evaluates
$\LEdiv(x, \EML(y,1))$ at all 8 witness points in $\mathcal{W}$.
Results logged in \texttt{python/results/s105\_door5\_sub\_2node.json}:
\[
  \text{all\_passed} = \texttt{True},\quad \text{max\_error} = 0.0.
\]

\noindent\textbf{Step 5: Optimality.}
By Theorem~\ref{thm:lower1}, no one-node tree in $\mathcal{F}_{16}$ computes
$x-y$.  Hence the two-node construction is optimal.
\end{proof}

\begin{remark}[Algebraic key identity]\label{rem:key}
The construction exploits the chain
\[
  \ln\!\bigl(e^x / e^y\bigr) = \ln\!\bigl(e^{x-y}\bigr) = x - y,
\]
which is the logarithm of a ratio of exponentials.
The operator $\LEdiv(a,b) = \ln(e^a/b)$ provides the ``ln-of-ratio'' in one
node; the inner node $\EML(y,1)=e^y$ converts $y$ to its exponential form,
making the denominator $e^y$ rather than a raw variable.
This is structurally analogous to the $\mul$ construction
$\ELAd(\EXL(0,x),y) = xy$ (T10u), where $\EXL(0,x) = \ln x$ converts $x$
to its logarithmic form.
\end{remark}

\subsection{The complete set of 2-node solutions}

\begin{proposition}[Complete 2-node solution set for sub in $\mathcal{F}_{16}$]\label{prop:six}
The exhaustive search over all 2-node trees in $\mathcal{F}_{16}$ with
terminal set $\mathcal{T}$ finds exactly \textbf{six} trees computing $x-y$:
\begin{align*}
  S_1 &= \LEdiv\!\bigl(x,\;\EML(y,1)\bigr),  \\
  S_2 &= \LEdiv\!\bigl(x,\;\EAL(y,1)\bigr),  \\
  S_3 &= \LEprod\!\bigl(x,\;\DEML(y,1)\bigr), \\
  S_4 &= \LEprod\!\bigl(x,\;\DEAL(y,1)\bigr), \\
  S_5 &= \LEdiv\!\bigl(x,\;\ELAd(y,1)\bigr),  \\
  S_6 &= \LEdiv\!\bigl(x,\;\ELSb(y,1)\bigr).
\end{align*}
\end{proposition}

\begin{proof}
These are the direct output of the certified exhaustive search logged in
\texttt{python/results/s105\_door5\_sub\_2node.json} (field
\texttt{full\_2node\_search.found}).
We verify correctness for each:

\medskip
\noindent$S_1$: $\EML(y,1)=e^y-0=e^y$;\quad
$\LEdiv(x,e^y)=\ln(e^x/e^y)=x-y$. \checkmark

\medskip
\noindent$S_2$: $\EAL(y,1)=e^y+\ln 1=e^y$;\quad
$\LEdiv(x,e^y)=x-y$. \checkmark
\smallskip
(EAL with second argument 1 degenerates to $e^y$, same as $S_1$.)

\medskip
\noindent$S_3$: $\DEML(y,1)=e^{-y}-\ln 1=e^{-y}$;\quad
$\LEprod(x,e^{-y})=\ln(e^x\cdot e^{-y})=\ln(e^{x-y})=x-y$. \checkmark

\medskip
\noindent$S_4$: $\DEAL(y,1)=e^{-y}+\ln 1=e^{-y}$;\quad
$\LEprod(x,e^{-y})=x-y$. \checkmark

\medskip
\noindent$S_5$: $\ELAd(y,1)=e^y\cdot 1=e^y$;\quad
$\LEdiv(x,e^y)=x-y$. \checkmark

\medskip
\noindent$S_6$: $\ELSb(y,1)=e^y/1=e^y$;\quad
$\LEdiv(x,e^y)=x-y$. \checkmark

\medskip
All six constructions collapse to the same algebraic chain: the inner node
produces $e^y$ (via different operators, all yielding $e^y$ when the constant
argument is 1), and the outer node applies $\LEdiv$ or $\LEprod$ to compute
$x-y$.  The six trees form two symmetric families:
$\{S_1,S_2,S_5,S_6\}$ using $\LEdiv$ with $e^y$ in the denominator, and
$\{S_3,S_4\}$ using $\LEprod$ with $e^{-y}$ as a multiplicative factor.
\end{proof}

% ════════════════════════════════════════════════════════════════════════════
\section{Structural Analysis}
% ════════════════════════════════════════════════════════════════════════════

\begin{remark}[Why the original F6 requires 3 nodes for sub]\label{rem:F6}
The classic three-node construction for $x-y$ over $\mathcal{F}_6$ is:
\[
  \EML\!\bigl(\EXL(0,x),\;\EML(y,1)\bigr)
  = e^{\ln x} - \ln(e^y - 0)
  = x - y.
\]
This uses 3 nodes: $\EXL(0,x)=\ln x$, $\EML(y,1)=e^y$, and the outer
$\EML(\ln x, e^y) = e^{\ln x} - \ln(e^y) = x - y$.

The obstruction in $\mathcal{F}_6$ is that every operator has the form
$e^{\pm a}\star\ln b$, which \emph{mandates} a logarithm on the second
argument.  To compute $x-y$ via $\ln(e^x/e^y)$, one needs:
\begin{enumerate}
  \item a node that produces $e^y$ (e.g., $\EML(y,1)$),
  \item a node that computes $e^x / e^y$, but all $\mathcal{F}_6$ operators
        apply $\ln$ to their second argument, so $e^y$ enters as $\ln(e^y)=y$
        rather than as $e^y$ directly, and
  \item a further node to recover $x-y$ from the available sub-expressions.
\end{enumerate}

The operator $\LEdiv(a,b) = \ln(e^a/b)$ is in $\mathcal{F}_{16}\setminus\mathcal{F}_6$:
it accepts $b$ \emph{directly} (not as $\ln b$), so $\LEdiv(x, e^y)$ does not
apply a second logarithm to $e^y$.  This is the structural reason why
$\mathcal{F}_{16}$ achieves 2 nodes while $\mathcal{F}_6$ requires 3.
\end{remark}

\begin{remark}[Structural parallel with mul]\label{rem:parallel}
The pair of results T10u (mul in 2 nodes) and T33 (sub in 2 nodes) share a
common structural pattern:

\smallskip
\begin{tabular}{@{}lll@{}}
\toprule
Operation & Inner node & Outer node \\
\midrule
$xy$ (T10u) & $\EXL(0,x) = \ln x$ & $\ELAd(\ln x, y) = e^{\ln x}\cdot y = xy$ \\
$x-y$ (T33) & $\EML(y,1) = e^y$ & $\LEdiv(x, e^y) = \ln(e^x/e^y) = x-y$ \\
\bottomrule
\end{tabular}
\smallskip

In both cases, the inner node converts one variable to a ``compatible form''
($\ln x$ for mul, $e^y$ for sub), and the outer node combines it with the
other variable via a semidirect-product-style operation.
The key operators $\ELAd$ and $\LEdiv$ are both members of $\mathcal{F}_{16}
\setminus\mathcal{F}_6$, specifically the operators that accept one argument
\emph{directly} without applying exp or ln to it.
\end{remark}

% ════════════════════════════════════════════════════════════════════════════
\section{Updated SuperBEST Table (v3)}
% ════════════════════════════════════════════════════════════════════════════

\subsection{SuperBEST v3: 19 nodes, 74.0\% savings}

\begin{theorem}[SuperBEST v3 is optimal]\label{thm:superbest3}
The minimum total node count for computing all seven standard arithmetic
operations over $\mathcal{F}_{16}$ with terminal set $\{x,y,0,1\}$ is
exactly \textbf{19 nodes}, yielding compression savings of
\[
  1 - \frac{19}{73} \;=\; \frac{54}{73} \;\approx\; 73.97\% \;\approx\; \mathbf{74.0\%}
\]
over the naive EML baseline of 73 nodes.
\end{theorem}

\begin{proof}
We argue row by row; the full table is Table~\ref{tab:superbest3}.

\noindent\textbf{exp(x), 1 node:}
$\EML(x,1)=e^x$ is a 1-node construction; 0 nodes is impossible
(no terminal equals $e^x$).  Optimum: 1.

\noindent\textbf{ln(x), 1 node:}
$\EXL(0,x)=\ln x$ is a 1-node construction; 0 nodes impossible.  Optimum: 1.

\noindent\textbf{neg(x), 2 nodes:}
T09 certifies a 2-node construction.  Exhaustive 1-node search confirms no
solution (T09 proof, certified in \texttt{python/results/s102\_door1\_mul\_1node.json}).
Optimum: 2.

\noindent\textbf{recip(x), 2 nodes:}
EDL construction is 2 nodes; 1-node lower bound by exhaustive search.
Optimum: 2.

\noindent\textbf{mul(x,y), 2 nodes:}
T10u certifies 2 nodes; 1-node lower bound by exhaustive search.  Optimum: 2.

\noindent\textbf{sub(x,y), 2 nodes:}
T33 (Theorem~\ref{thm:T33}) certifies 2 nodes.  1-node lower bound is
Theorem~\ref{thm:lower1}.  Optimum: 2.

\noindent\textbf{add(x,y), 3 nodes:}
T11 certifies 3 nodes.  Exhaustive 2-node search over $\mathcal{F}_{16}$
finds zero solutions.  Optimum: 3.

\noindent\textbf{pow(x,n), 3 nodes:}
T11 construction via EXL certifies 3 nodes.  Exhaustive 2-node search
finds zero solutions.  Optimum: 3.

\smallskip
Summing per-row optima: $1+1+2+2+2+2+3+3 = \mathbf{19}$.
The row-by-row lower bounds are independent, so the global minimum is 19.
The savings ratio is $1 - 19/73 = 54/73 \approx 0.7397 \approx 74.0\%$.
\end{proof}

\begin{table}[h]
\centering
\caption{SuperBEST v3: Minimum node counts and savings.
         The sub row (highlighted) is updated by T33.
         The mul row was updated by T10u (SuperBEST v2).
         All other rows are unchanged from SuperBEST v1 (T08).}
\label{tab:superbest3}
\smallskip
\begin{tabular}{@{}lccccl@{}}
\toprule
\textbf{Operation} & \textbf{Naive-EML} & \textbf{Best-F6} &
\textbf{Best-F16} & \textbf{Savings (F16 vs Naive)} & \textbf{Reference} \\
\midrule
$\operatorname{exp}(x)=e^x$        & 1  & 1 & 1 & $0\%$             & trivial \\
$\operatorname{ln}(x)=\ln x$       & 3  & 1 & 1 & $66.7\%$          & T01 \\
$\operatorname{neg}(x)=-x$         & 7  & 2 & 2 & $71.4\%$          & T09 \\
$\operatorname{recip}(x)=1/x$      & 5  & 2 & 2 & $60.0\%$          & T09 \\
$\operatorname{mul}(x,y)=xy$       & 13 & 3 & 2 & $84.6\%$          & T10u \\
$\operatorname{sub}(x,y)=x-y$      & 11 & 3 & \textbf{2} & $\mathbf{81.8\%}$ & \textbf{T33} \\
$\operatorname{add}(x,y)=x+y$      & 13 & 3 & 3 & $76.9\%$          & T11 \\
$\operatorname{pow}(x,n)=x^n$      & 15 & 3 & 3 & $80.0\%$          & EXL \\
\midrule
\textbf{Total}                     & \textbf{73} & \textbf{21} & \textbf{19} &
  $\mathbf{74.0\%}$ & \\
\bottomrule
\end{tabular}

\smallskip
\footnotesize{%
\textit{Notes.}
(1) ``Naive-EML'' counts nodes in a direct EML-only tree for each operation
without optimization.
(2) Best-F6 is the minimum node count using operators from $\mathcal{F}_6$ only;
Best-F16 uses the full extended family $\mathcal{F}_{16}$.
(3) Savings: $(1 - \text{Best-F16}/\text{Naive-EML})$.
(4) The total savings of $74.0\% = 1 - 19/73 = 54/73$ is the second consecutive
update to the SuperBEST table within the same sprint (after T10u raised it from
71.2\% to 72.6\%).
(5) This is the \emph{second consecutive 2-node result}: both mul and sub,
previously believed to require 3 nodes in $\mathcal{F}_{16}$, admit 2-node
constructions using operators from $\mathcal{F}_{16}\setminus\mathcal{F}_6$.
}
\end{table}

% ════════════════════════════════════════════════════════════════════════════
\section*{Computational Artifacts}
% ════════════════════════════════════════════════════════════════════════════

All search results are reproducible via:
\begin{verbatim}
python python/scripts/door5_sub_2node.py
\end{verbatim}
Output: \texttt{python/results/s105\_door5\_sub\_2node.json}.

Key fields:
\begin{verbatim}
{
  "upper_bound": { "all_passed": true, ... },
  "lower_bound_1node": {
    "total_1node_trees_checked": 256,
    "matches_found": 0,
    "lower_bound_holds": true
  },
  "full_2node_search": {
    "count": 6,
    "found": [
      {"tree": "LEdiv(x,EML(y,1))"},
      {"tree": "LEdiv(x,EAL(y,1))"},
      {"tree": "LEprod(x,DEML(y,1))"},
      {"tree": "LEprod(x,DEAL(y,1))"},
      {"tree": "LEdiv(x,ELAd(y,1))"},
      {"tree": "LEdiv(x,ELSb(y,1))"}
    ]
  },
  "conclusion": { "minimum_nodes": 2, "T33_confirmed": true }
}
\end{verbatim}

% ════════════════════════════════════════════════════════════════════════════
\section*{Theorem Catalog Labels}
% ════════════════════════════════════════════════════════════════════════════

\begin{description}
  \item[T33a]  No 1-node tree in $\mathcal{F}_{16}$ with terminals $\{x,y,0,1\}$
               computes $x-y$; 256 candidates exhaustively searched, zero matches
               (Theorem~\ref{thm:lower1}, this paper).
  \item[T33]   $\sub(x,y) = \LEdiv(\EML(y,1), x)$; 2-node optimal in $\mathcal{F}_{16}$;
               six 2-node solutions found by exhaustive search
               (Theorem~\ref{thm:T33}, this paper).
  \item[SuperBEST v3]  Global minimum 19 nodes over $\mathcal{F}_{16}$;
               74.0\% savings over naive EML baseline of 73 nodes
               (Theorem~\ref{thm:superbest3}, this paper).
\end{description}

% ════════════════════════════════════════════════════════════════════════════
\section*{Acknowledgments}
% ════════════════════════════════════════════════════════════════════════════

This work continues the monogate SuperBEST sprint (sessions S100--S105).
T33 was discovered during session S105 after the primary algebraic identity
$\LEdiv(x, e^y) = x - y$ was verified computationally.

\bigskip
\noindent\textit{Almaguer, A.R.\ (2026).
Subtraction in Two Nodes: T33 and the SuperBEST v3 Update to 19 Nodes.
monogate research. \url{https://monogate.org}}

\end{document}
